\title{\thetitle}
\author{\theauthors}
\date{}
\maketitle

\begin{abstract}
Мы исследуем, как квадратичные ограничения на шесть расстояний между точками каждой четвёрки влияют на геометрию пространств с внутренней метрикой.
\end{abstract}

\section{Введение}\label{par:quadratic-inq}

\paragraph{Квадратичное условие.}
Пусть $\bm{x}\z=(x_1,\ldots,x_n)$ --- набор из $n$ точек метрического пространства $X$, а $a_{i,j}$ --- элементы симметричной матрицы размера $n{\times}n$.
Неравенство вида
\[\sum_{i,j}a_{i,j}\cdot|x_i-x_j|_X^2\geqslant 0\]
будет называться \emph{квадратичным}, а система квадратичных неравенств --- \emph{квадратичным условием}.

Нас будут интересовать пространства с внутренней метрикой, в которых заданное квадратичное условие выполнено для любой $n$-ки точек.
Следующее утверждение говорит, что у теоремы Топоногова нет простых аналогов.
Оно будет доказано в разделе~\ref{par:globalization}.

\begin{thm}{Основная теорема}
Предположим, что квадратичное условие на 4-ки точек обладает свойством глобализации;
то есть если в пространстве с внутренней метрикой это условие выполнено локально (в окрестности каждой точки), то оно выполнено глобально.
Тогда либо условие тривиально, либо оно задаёт класс александровских пространств неотрицательной кривизны.
\end{thm}

\paragraph{Вспомогательные результаты.}
Сформулируем пару вспомогательных результата, представляющих самостоятельный интерес.
В их формулировках используется частный случай квадратичных неравенств
с $a_{i,j}=-\lambda_i\cdot\lambda_j$, где $(\lambda_1,\ldots, \lambda_n)$ --- набор вещественных чисел, такой что $\lambda_1+\ldots+\lambda_n=0$.
Это так называемые неравенства отрицателного типа, они обсудаются в разделе~\ref{par:rank-one}.

В разделе~\ref{Four-point arrays} мы докажем следующий вариант теоремы Абрахама Вальда \cite[§ 7]{wald}, описывающий метрики на всевозможных четвёрках точек в александровских пространствах неотрицательной и неположительной кривизны.

\begin{thm}{Предложение}
Следующие три условия эквивалентны:
\begin{enumerate}[(i)]
\item Метрическое пространство $X$, состоящее из четырёх точек, изометрично подмножеству пространства с внутренней метрикой и неотрицательной (неположительной) кривизной в смысле Александрова.
\item $X$ изометрично подмножеству произведения $\RR^3\times r\cdot \SSS^1$ для некоторого $r>0$ (соответственно, произведения $\RR^3\times Y$, где $Y$ обозначает треногу, то есть три полупрямые с общим началом).
\item В пространстве $X$ выполнены все неравенства отрицательного типа с $\lambda$-наборами, для которых $\lambda_1\cdot\lambda_2\cdot\lambda_3\cdot\lambda_4<0$ (соответственно, с $\lambda$-наборами, для которых $\lambda_1\cdot\lambda_2\cdot\lambda_3\cdot\lambda_4>0$).
\end{enumerate}
\end{thm}

Седующее предложение говорит, что каждое невырожденное неравенство отрицательного типа задаёт класс александровских пространств либо неотрицательной, либо неположительной кривизны; оно будет доказано в разделе~\ref{Alexandrov's comparison}.

\begin{thm}{Предложение}
Рассмотрим класс всех пространств с внутренней метрикой, удовлетворяющих квадратичному условию, заданному неравенством отрицательного типа с $\lambda$-набором $(\lambda_1,\lambda_2,\lambda_3,\lambda_4)$.
\begin{enumerate}[(i)]
\item Если $\lambda_1\cdot\lambda_2\cdot\lambda_3\cdot\lambda_4<0$, то это условие задаёт класс всех александровских пространств неотрицательной кривизны.
\item Если $\lambda_1\cdot\lambda_2\cdot\lambda_3\cdot\lambda_4>0$, то это условие задаёт класс всех александровских пространств неположительной кривизны.
\item Если $\lambda_1\cdot\lambda_2\cdot\lambda_3\cdot\lambda_4=0$, то это условие не накладывает никаких ограничений; ему удовлетворяют все пространства с внутренней метрикой.
\end{enumerate}

\end{thm}

Это утверждение показывает, что сравнительно слабая форма метрического сравнения влечёт значительно более сильную форму (при условии, что метрика внутренняя).
Доказательство использует рассуждение Такаси Сато \cite{sato} и его модификацию, предложенную двумя авторами \cite{lebedeva-petrunin-2010};
в этих статьях рассматриваются частные случаи таких неравенств для $\lambda$-наборов $(1,1,-1,-1)$ и $(1,1,1,-3)$.

В разделах \ref{Associated form} и \ref{par:rank-one} вводятся необходимые определения.
В них мы определяем квадратичную форму ассоциированную с набором точек и неравенства отрицательного типа.
Мы также доказываем несколько утверждений, справедливых для всех $n$-ок точек.
В разделе \ref{Auxiliary statements} приведено техническое утверждение, используемое в доказательстве основной теоремы.
